% !TEX root = main.tex
\subsection{Correctness and Security}\label{subsection:sms-correctness-and-security}

The system maintains a public state $\State$ recording each enrolled user's
attribute, assigned function, and enrollment transcript.
$\State$ is \emph{consistent} if every recorded transcript verifies
under $\SMSJoin$.
A secure \SMS scheme satisfies correctness, privacy, and unforgeability.

\medskip
\noindent\textbf{Correctness.}
Correctness requires that honest signatures verify and open correctly,
that the system state remains consistent throughout,
and that only enrolled users can produce valid signatures.

\begin{definition}[Correctness]\label{definition:sms-correctness}
For $\SMSSetup(1^\lambda)$ returning $(\pp, \State, \mfk, \tsk)$,
an \SMS scheme $\SMSScheme$ is \emph{correct} if the following conditions hold.
\begin{enumerate}
\item $\OpeningSetSignable \subsetneq \OpeningSet$
and all attributes in $\State$ are pairwise distinct.

\item If
$\SMSJoin\interaction{\User(1^\lambda)}{\FA(\mfk)}(\pp, \State, \attr, F)$
is executed honestly, with transcript $\transcript$:
$\certattr \!\leftrightharpoons_\pp\! \secattr$,
$\transcript$ is valid,
with $(\attr, F, \transcript)$ inserted into~$\State$.

\item For $(\attr, F, \cdot) \in \State$ with
$\certattr \leftrightharpoons_\pp \secattr$,
$G \in \GFamily$, $M \in \MessageSpace$, $w \in \WitnessSpace$:
\begin{equation}\label{equation:sms-correctness-probability}
\condprob{
\phantom{\land}\; b = 1
\\
\land\; \op \in \OpeningSetSignable
\\
\land\; \op = \op' 
\\
\land\; \tM = M'
}{
(\tM, \Sigma) \gets \SMSSign(\pp, \certattr, \secattr, G, M, w),
\\
(M', \op) = P(F(M, \attr, w), G(M, \attr, w)),
\\
b \coloneqq \SMSVerify(\pp, G, \tM, \Sigma),
\\
\op' \gets \SMSOpen(\pp, \tsk, G, \tM, \Sigma)
} = 1.
\end{equation}
\end{enumerate}
\end{definition}

\medskip
\noindent\textbf{Adversary model.}
The adversary $\Adversary$ interacts with the system via oracles
that model realistic attack capabilities,
defined below,
and the security properties are parameterized by which oracles $\Adversary$ may access.
We write $\!(\cdot)^{\condition}$ to denote that an oracle returns $\epsilon$
when $\condition$~holds.

\medskip
\noindent\underline{Oracles $\OMfk$, $\OTsk$.}
$\OMfk$ returns $\mfk$ to $\Adversary$ and sets flag $\amfk$ to $\True$.
$\OTsk$ returns $\tsk$ and sets flag $\atsk$ to $\True$.
Both flags are initialized to $\False$.

\medskip
\noindent\underline{Oracle $\OCorruptU$, set $\LCU$.}
On query $(\attr, F)$ with $(\attr, \cdot) \notin \LCU$,
$\OCorruptU$ acts as an honest $\FA$, runs $\SMSJoin$,
sets $\LCU \coloneqq \LCU \cup \{(\attr, F)\}$,
updates $\State$, and returns $(\secattr, \certattr)$ to $\Adversary$.

\medskip
\noindent\underline{Oracle $\OHonestU$, set $\LHU$.}
On query $\attr$ with $(\attr, \cdot) \notin \LCU$,
$\OHonestU$ acts as an honest $\FA$, samples $F \in \FFamily$, runs $\SMSJoin$,
and sets $\LHU \coloneqq \LHU \cup \{(\attr, F)\}$ and updates $\State$.
The secrets of users in $\LHU$ are not given to $\Adversary$.
Both $\OCorruptU$ and $\OHonestU$ guarantee that
all attributes in $\State$ are distinct
and $\LCU \cap \LHU = \emptyset$ on the attribute component.

\medskip
\noindent\underline{Oracle $\OSign$, set $\LSig$.}
On query $(\attr, G, M, w)$,
if $\exists F$ such that $(\attr, F) \in \LHU$,
$G \in \GFamily$, $M \in \MessageSpace$, $w \in \WitnessSpace$,
and $\op \in \OpeningSetSignable$ where $\op$ is determined by $P$, $F$, $G$, $M$, $\attr$, $w$,
then $\OSign$ returns
$(\tM, \Sigma) \gets \SMSSign(\pp, \certattr, \secattr, G, M, w)$
and inserts $(\attr, G, M, \Sigma)$ into $\LSig$.

\medskip
\noindent\underline{Oracle $\OOpen$.}
On query $(G, \tM, \Sigma)$,
return $\op \gets \SMSOpen(\pp, \tsk, G, \tM, \Sigma)$.

\medskip
\noindent\textbf{Privacy.}
Privacy guarantees that \SMS signatures reveal no information about
the signer's attribute $\attr$ or hidden policy $F$.
$\Adversary$ chooses $G$ and two candidate signers,
each with a valid certificate, secret, message, and witness,
and attempts to identify which signer produced the challenge signature.
Both candidate signers must use the same $G$,
since $G$ is public and a verifier could trivially distinguish
signatures produced under different public functions.

$\Adversary$ may access $\OCorruptU$, $\OHonestU$, $\OMfk$, and $\OSign$,
but not $\OTsk$ or $\OOpen$:
access to $\tsk$ trivially breaks privacy,
and allowing $\OOpen$ access requires public-key encryption to be secure against chosen-ciphertext attacks,
a conceptually straightforward extension (\eg,~\cite{CRYPTO:NGSY22}).
To prevent trivial wins, both candidate signers must yield
the same derived message $\tM$ and a signable opening value.

\begin{definition}[Privacy]\label{definition:sms-privacy}
Let $\ExpPrivb_{\Adversary}(\lambda)$ be defined as in
Figure~\ref{figure:sms-security-experiments}.
$\SMSScheme$ is \emph{private} if, for any probabilistic polynomial-time (PPT) $\Adversary$,
\[
\left|\Pr\bigl[\ExpPrivZero_{\Adversary}(\lambda) = 1\bigr]
- \Pr\bigl[\ExpPrivOne_{\Adversary}(\lambda) = 1\bigr]\right|
\leq \negl(\lambda).
\]
\end{definition}

% !TEX root = main.tex
\begin{experiment}[Experiments for security of \SMS\label{figure:sms-security-experiments}]

\underline{\textbf{Experiment $\ExpPrivb_{\Adversary}(\lambda)$:}}

$(\pp, \State, \mfk, \tsk) \gets \SMSSetup(1^\lambda)$.

$\amfk \coloneqq \False$,
$\LCU \coloneqq \emptyset$,
$\LHU \coloneqq \emptyset$,
$\LSig \coloneqq \emptyset$.

$(\aux, G, (\certattri, \secattri, M_i, w_i)_{i \in \{0, 1\}})
\gets \Adversary^{\OCorruptU, \OHonestU, \OMfk, \OSign}(\pp)$.

Obtain $(\attr_i, F_i)$ from $\certattri$ for $i \in \{0, 1\}$.

If $G \notin \GFamily$ or $(\exists i \in \{0, 1\} \colon M_i \notin \MessageSpace \lor \attr_i \notin \AttrSpace
\lor F_i \notin \FFamily
\lor \secattri \not\leftrightharpoons_\pp \certattri)$,
return $0$.

$(\tM_i, \op_i) \coloneqq P(F_i(M_i, \attr_i, w_i), G(M_i, \attr_i, w_i))$
for $i \in \{0, 1\}$.

If $\tM_0 \neq \tM_1
\lor \op_0 \notin \OpeningSetSignable
\lor \op_1 \notin \OpeningSetSignable$,
return $0$.

Set $(\tM, \Sigma) \gets \SMSSign(\pp, \certattrb, \secattrb, G, M_b, w_b)$.

$b' \gets \Adversary^{\OCorruptU, \OHonestU, \OMfk, \OSign}(\pp, \tM, \Sigma, \aux)$.

Return $b'$.

\vspace{1ex}
\hrule height 0.4pt
  
\medskip
\underline{\textbf{Experiment $\ExpReal_{\Adversary}(\lambda)$:}}

$(\pp, \State, \mfk, \tsk) \gets \SMSSetup(1^\lambda)$.

$\amfk \coloneqq \False$, $\atsk \coloneqq \False$,
$\LCU \coloneqq \emptyset$, $\LHU \coloneqq \emptyset$, $\LSig \coloneqq \emptyset$.

$b \gets \Adversary^{\OCorruptU, \OHonestU, \OMfk, \OTsk, \OSign, \OOpen}(\pp)$.

Return $b$.

\vspace{1ex}
\hrule height 0.4pt

\medskip
\underline{\textbf{Experiment $\ExpSim_{\Adversary}(\lambda)$:}}

$(\pp, \State, \mfk, \tsk, \tauext) \gets \SMSSimSetup(1^\lambda)$.
\cmmt{Change $\SMSSetup(1^\lambda)$ in $\ExpReal_{\Adversary}(\lambda)$
to $\SMSSimSetup(1^\lambda)$\!}

$\amfk \coloneqq \False$, $\atsk \coloneqq \False$,
$\LCU \coloneqq \emptyset$, $\LHU \coloneqq \emptyset$, $\LSig \coloneqq \emptyset$.

$b \gets \Adversary^{\OCorruptU, \OHonestU, \OMfk, \OTsk, \OSign, \OOpen}(\pp)$.

Return $b$.

\vspace{1ex}
\hrule height 0.4pt

\medskip
\underline{\textbf{Experiment $\ExpUnforge_{\Adversary}(\lambda)$:}}

$(\pp, \State, \mfk, \tsk, \tauext) \gets \SMSSimSetup(1^\lambda)$.

$\amfk \coloneqq \False$, $\atsk \coloneqq \False$,
$\LCU \coloneqq \emptyset$, $\LHU \coloneqq \emptyset$, $\LSig \coloneqq \emptyset$.

$(G, \tM, \Sigma) \gets
\Adversary^{\OCorruptU, \OHonestU, \OMfk, \OTsk, \OSign, \OOpen}(\pp)$.

If $G \notin \GFamily
\lor \tM \notin \TMessageSpace
\lor \SMSVerify(\pp, G, \tM, \Sigma) = 0$,
return $0$.

$(\attr', F', M', w') \gets \SMSExtract(\pp, \tauext, G, \tM, \Sigma)$.

If $\attr' \notin \AttrSpace
\lor F' \notin \FFamily
\lor M' \notin \MessageSpace
\lor w' \notin \WitnessSpace
\lor (\attr', G, M', \Sigma) \in \LSig$,
return $0$.

$\op \gets \SMSOpen(\pp, \tsk, G, \tM, \Sigma)$.

$(\tildeM, \op') \coloneqq P(F'(M', \attr', w'), G(M', \attr', w'))$.

Return $1$ if \eqref{equation:condition-for-winning-unforgeability} holds;
otherwise return $0$.
\begin{equation}\label{equation:condition-for-winning-unforgeability}
\begin{aligned}
& (\op' \notin \OpeningSetSignable
	\lor \op \neq \op'
	\lor \tM \neq \tildeM)
	&& \cmmt{extractability}
	\\
& \lor (\amfk = \False \land (\attr', F') \notin \LCU)
	&& \cmmt{type-1 unforgeability}
	\\
& \lor (\amfk = \False \land (\attr', F') \in \LHU)
	&& \cmmt{type-2 unforgeability}
\end{aligned}
\end{equation}

\end{experiment}


\medskip
\noindent\textbf{Unforgeability.}
Unforgeability prevents any PPT $\Adversary$ from producing a valid \SMS signature
without a legitimate certificate--secret pair.
It encompasses three nested properties,
formalized jointly in $\ExpUnforge_\Adversary(\lambda)$
(Figure~\ref{figure:sms-security-experiments}).

\begin{itemize}
\item \emph{Extractability.}
Every valid forgery $(G, \tM, \Sigma)$ must extract to a list of witnesses
$(\attr', F', M', w')$ such that
$\op' \in \OpeningSetSignable$,
$\op' = \op$ as opened by $\SMSOpen$,
and $\tildeM = \tM$.
No adversary can produce a signature whose extracted witness is inconsistent
with its opening value or its derived message.

\item \emph{Type-1 unforgeability (traceability).}
Extractability holds, and $\Adversary$ cannot forge a signature
whose extracted $(\attr', F')$ lies outside the corrupted set $\LCU$.

\item \emph{Type-2 unforgeability (non-frameability).}
Extractability holds, and no coalition of users, $\FA$, and $\TA$
can forge a signature whose extracted $(\attr', F')$ belongs to $\LHU$.
\end{itemize}

The two unforgeability types are independent, as the table below shows.

\begin{center}\footnotesize
\begin{tabular}{@{} l | l | l @{}}
\hline
& \textbf{Type-1 (Traceability)} & \textbf{Type-2 (Non-frameability)}
\\
\hline
Forgery target
	& $(\attr', F') \notin \LCU$
	& $(\attr', F') \in \LHU$
	\\
Covers
	& $\LHU \cup (\AttrSpace \setminus (\LCU \cup \LHU))$
	& $\LCU \cup (\AttrSpace \setminus (\LCU \cup \LHU))$
	\\
$\OMfk$ access
	& Restricted & Restricted
	\\
\hline
\end{tabular}
\end{center}

Neither type implies the other:
Type-2 covers $\LCU$ but not $\LHU$;
Type-1 covers $\LHU$ but not $\LCU$.
A forgery is non-trivial only if $(G, \tM, \Sigma)$ was not previously
obtained via $\OSign$, enforced by the $\LSig$ membership check in $\ExpUnforge$.

Extractability permits $\OMfk$ access in $\ExpUnforge$ because it only requires consistency of extracted components with the system prescription, independent of membership.
Conversely, for Type-1 and \mbox{Type-2} unforgeability, $\OMfk$ access is restricted,
because possessing $\mfk$ would allow $\Adversary$ to certify attributes outside both $\LCU$ and $\LHU$, trivially breaking the membership-based winning conditions;
in particular, $\Adversary$ could certify the same attribute as an honest user
under a different hidden function
(\eg, $(\attr, F) \in \LHU$ but $(\attr, F') \notin \LHU$ where $F' \neq F$),
producing a valid signature outside the honest-user record.

\medskip
\noindent\textbf{Simulated setup for extraction.}
%Following~\cite{PKC:BelFuc14,EC:LNPY21,CRYPTO:NGSY22,AC:NRSX23,ACNS:HLTW25},
We devise algorithms $\SMSSimSetup$ and $\SMSExtract$
to enable simulation-sound extractability~\cite{CRYPTO:ChaLys06}:
a simulated setup indistinguishable from the real one,
%paired
with a trapdoor for extracting the underlying~witness.
%any valid signature.

\begin{definition}[Simulated setup and extraction]%
\label{definition:sms-additional-algorithms}
\begin{description}
\item[$\SMSSimSetup(1^\lambda)$:]
Identical to $\SMSSetup$ but additionally outputs extraction trapdoor $\tauext$.
Output: $(\pp, \State, \mfk, \tsk, \tauext)$.

\item[$\SMSExtract(\pp, \tauext, G, \tM, \Sigma)$:]
Given $\pp$, trapdoor $\tauext$, $G \in \GFamily$,
$\tM \in \TMessageSpace$, and signature $\Sigma$,
return $(\attr, F, M, w) \in
\AttrSpace \times \FFamily \times \MessageSpace \times \WitnessSpace$.
\end{description}
\end{definition}

The correctness of \SMS under $\SMSSimSetup$ is defined identically
to Definition~\ref{definition:sms-correctness} with $\SMSSetup$
replaced by $\SMSSimSetup$;
additionally, the extracted components must match those used to create the signature.
Setup indistinguishability implies this correctness against PPT adversaries,
while correctness of the simulated setup holds unconditionally.

\begin{definition}[Setup indistinguishability]%
\label{definition:sms-setup-indistinguishability}
$\SMSScheme$ satisfies \emph{setup indistinguishability} if,
for any PPT $\Adversary$,
\[
\left|\Pr\bigl[\ExpReal_{\Adversary}(\lambda) = 1\bigr]
- \Pr\bigl[\ExpSim_{\Adversary}(\lambda) = 1\bigr]\right|
\leq \negl(\lambda),
\]
where $\ExpReal_{\Adversary}(\lambda)$ and $\ExpSim_{\Adversary}(\lambda)$
are defined in Figure~\ref{figure:sms-security-experiments}.
\end{definition}

\begin{definition}[Unforgeability]\label{definition:sms-unforgeability}
$\SMSScheme$ is \emph{unforgeable} if, for any PPT $\Adversary$,
\[
\Pr\bigl[\ExpUnforge_{\Adversary}(\lambda) = 1\bigr] \leq \negl(\lambda),
\]
where $\ExpUnforge_{\Adversary}(\lambda)$ is defined in
Figure~\ref{figure:sms-security-experiments}.
\end{definition}
